Self-critique and iterative refinement are increasingly used to improve large language model (LLM) outputs, yet it remains unclear whether these techniques induce genuine restructuring of internal representations or merely produce shallow re-sampling. We investigate this question by extracting residual stream activations across all 28 layers of \qwen at three pipeline stages---initial answer, self-critique, and revised answer---plus a control condition (re-prompting without critique), on 200 \gsmk math problems. We find that self-critique induces significantly larger representational drift than simple re-prompting (mean cosine similarity 0.57 vs.\ 0.69, $p < 0.01$ at all 28 layers after FDR correction), confirming that critique produces genuine internal restructuring beyond prompt variation. However, this restructuring is \emph{destructive}: accuracy drops from 90.5\% to 84.0\% after self-critique, while re-prompting achieves 92.0\%. Linear probes trained to predict answer correctness show decreased separability in post-critique representations, and PCA analysis reveals that critique preserves the dominant variance direction but disrupts finer-grained features. Our results provide the first mechanistic evidence that self-critique acts as a perturbation that displaces representations from correct reasoning subspaces rather than steering them toward error correction, with implications for the design of self-improvement agent architectures.
